% Volume I: The Epistemology of Suspicion
% Part I. The Power of Paranoia
% Chapter 5: Distrust Without Accusation
% Spine: proof-spines/ch005-spine.md

\chapter{Distrust Without Accusation}
\label{ch:ch005}


Distrust need not begin in accusation. Institutions can be designed around fallibility while remaining agnostic about which failure, if any, has occurred. A protocol may assume
\[
P(\text{error})>0,\qquad
P(\text{capture})>0,\qquad
P(\text{misrepresentation})>0
\]
without claiming in advance that any particular actor is malicious. The point is to structure action under unresolved uncertainty rather than to smuggle in a hidden verdict of guilt.

On that basis, institutional distrust is expressed through redundancy, separation of powers, audit, cross-examination, reproducibility, conflict-of-interest rules, and appeal. These mechanisms encode a general proposition: no observer, witness, office, or archive should be treated as a transparent channel to reality. Distrust is therefore not identical with condemnation. It is a way of distributing vulnerability when the evidence does not yet justify accusation but does justify procedural safeguards.
